% Source: https://artofproblemsolving.com/community/c3h2615498_peace09s_3000th_post
% Stable identity: aops-topic-2615498-q26
Let $a=BC$, $b=CA$, $c=AB$, so $2a=b+c$. By the angle bisector theorem $BD=\dfrac{ac}{b+c}$.
\begin{asy}
usepackage("amsmath");
size(4cm);
defaultpen(fontsize(10pt));
pair B=(0,0),C=(5,0);pair A=(1.1,3.84);real a=abs(B-C),bb=abs(C-A),c=abs(A-B);pair I=(a*A+bb*B+c*C)/(a+bb+c);pair D=extension(A,I,B,C);
draw(A--B--C--cycle);draw(A--D,dashed);draw(B--I,dotted);dot(I);
label("$A$",A,N);label("$B$",B,SW);label("$C$",C,SE);label("$D$",D,S);label("$I$",I,E);
\end{asy}

In triangle $ABD$, $BI$ bisects $\angle B$, so again by the angle bisector theorem
\[\dfrac{AI}{ID}=\dfrac{AB}{BD}=\dfrac{c(b+c)}{ac}=\dfrac{b+c}{a}=2.\ \blacksquare\]
